\chapter*{ABSTRAK}
\addcontentsline{toc}{chapter}{ABSTRAK}

\begin{center}

    \textbf{Analisis dan Evaluasi Prediksi Harian PM10 Jakarta Menggunakan Model Hibrida Random Forest-ARIMA Berbasis Rekayasa Fitur} \\
    \vspace{0.4cm}
    oleh \\
    \vspace{0.4cm}
    Muhammad Rafi Dhiyaulhaq\\
    18222069\\
\end{center}

\begin{spacing}{1.0}
Penurunan kualitas udara di DKI Jakarta, khususnya akibat tingginya konsentrasi partikulat PM\textsubscript{10}, memerlukan sistem peringatan dini untuk membantu mitigasi dampak kesehatan masyarakat. Banyak model prediksi saat ini memproses data polutan mentah tanpa mengekstraksi komponen temporalnya, serta memiliki struktur yang sulit diinterpretasikan secara langsung oleh pemangku kebijakan. Selain itu, penggunaan model tunggal memiliki batasan operasional; \textit{Random Forest} dapat memetakan pola non-linier namun rentan terhadap autokorelasi pada residu, sedangkan model ARIMA efektif dalam memodelkan komponen linier deret waktu namun kurang adaptif terhadap lonjakan data yang ekstrem. Oleh karena itu, penelitian ini menerapkan arsitektur hibrida \textit{Random Forest}-ARIMA sekuensial yang dikombinasikan dengan teknik \textit{feature engineering}. Proses rekayasa fitur menghasilkan 15 variabel turunan, meliputi efek jeda waktu (\textit{lag}), rata-rata bergerak (\textit{rolling statistics}), pola temporal, dan interaksi antarpolutan, yang diekstrak dari dataset Indeks Standar Pencemar Udara (ISPU) DKI Jakarta periode 2010--2025. Evaluasi menggunakan \textit{time-series cross-validation} (5-\textit{fold}) menunjukkan bahwa model hibrida yang diusulkan menghasilkan nilai \textit{Root Mean Square Error} (RMSE) sebesar 13,4004. Nilai kesalahan prediksi tersebut lebih rendah jika dibandingkan dengan model ARIMA murni (RMSE 18,8432), \textit{Random Forest} murni (RMSE 13,5717), serta algoritma komparasi XGBoost (RMSE 14,3307). Hasil studi ablasi menunjukkan bahwa integrasi fitur turunan berkontribusi dalam menurunkan nilai galat prediksi. Model ini juga dilengkapi dengan metode \textit{SHapley Additive exPlanations} (SHAP) untuk menganalisis kontribusi setiap fitur terhadap hasil prediksi secara lebih transparan.

Kata kunci: PM\textsubscript{10}, Kualitas Udara, \textit{Random Forest}, ARIMA, \textit{Feature Engineering}, \textit{Explainable AI}, SHAP.
\end{spacing}